\section{The global Hartogs metric}\label{sec:hartogs}

We now turn to Proposition~\ref{metric:transfer}.  In this section we
prove all its global geometric assertions.  The real sectional-curvature
bounds will follow from the calculation in Section~\ref{sec:curvature}
and the estimates in Section~\ref{sec:pinching}.

The potential specializes Calabi's construction on a Hermitian line
bundle \cite[Section~3]{Calabi1979}. We verify its positivity,
smoothness, and completeness directly, including at the zero section.

\begin{lemma}\label{metric:complete}
Let $\phi$ be a smooth real plurisubharmonic function on $\B$, and let
$\lambda>0$.  The domain $M_\phi$ in \eqref{metric:domain} is
contractible, and the complex Hessian of
$\mathcal K_\lambda=\lambda\psi-\log(1-e^\phi|w|^2)$ defines a smooth
geodesically complete K\"ahler metric
$g=2\operatorname{Re}\widehat g$ on $M_\phi$.
Moreover, $\mathcal K_\lambda$ is a smooth strictly plurisubharmonic
exhaustion of $M_\phi$.
\end{lemma}

\begin{proof}
\emph{The domain and a smooth positive metric.}
The function $t=e^\phi|w|^2$ is smooth on $\B\times\C$, so
$M_\phi=\{t<1\}$ is open in $\C^3$.  The map
\[
 (z,w)\longmapsto (z,e^{\phi(z)/2}w)
\]
is a smooth real diffeomorphism from $M_\phi$ to $\B\times\mathbb D$,
where $\mathbb D=\{u\in\C:|u|<1\}$.  Its inverse is
$(z,u)\mapsto(z,e^{-\phi(z)/2}u)$.  Thus $M_\phi$ is contractible;
no holomorphic trivialization is needed for this conclusion.

The potential $\mathcal K_\lambda$ is smooth at every point of
$M_\phi$, including its zero section.  Write
\[
 h=\partial\bar\partial\phi,\qquad
 x=\frac{t}{1-t},\qquad G=\lambda I+xh,
 \qquad \vartheta=dw+w\partial\phi.
\]
The base covectors together with $\vartheta$ form a global smooth
complex coframe.  Differentiating the potential gives
\begin{equation}\label{metric:smooth-form}
 \widehat g
 =G+\frac{e^\phi}{(1-t)^2}\,
       \vartheta\otimes\overline{\vartheta}.
\end{equation}
For a complex tangent vector $u=(u_z,u_w)$ this means
\[
 \widehat g(u,u)
 =G(u_z,u_z)
  +\frac{e^\phi}{(1-t)^2}
      |u_w+w\partial\phi(u_z)|^2.
\]
Since $h\ge0$, the first term controls every nonzero horizontal
projection through $G\ge\lambda I$; when $u_z=0$, the second term is
positive unless $u_w=0$.  Thus $\widehat g$ is positive definite.
At the zero section its matrix is particularly simple:
\begin{equation}\label{metric:axis}
 \widehat g\big|_{(z,0)}
 =\lambda I(z)+e^{\phi(z)}\,dw\otimes d\bar w.
\end{equation}
The mixed blocks vanish.  Being a positive complex Hessian of a smooth
real potential, $\widehat g$ defines the stated K\"ahler metric.

Away from the zero section one can write $s=\log t$ and
$q=x(1+x)=t/(1-t)^2$.  Then $\partial s=\vartheta/w$, so
\eqref{metric:smooth-form} also reads
\begin{equation}\label{metric:log-form}
 \widehat g=\lambda I+xh+
             q\,\partial s\otimes\overline{\partial s}.
\end{equation}
We will use this expression for the curvature computation.  Formula
\eqref{metric:smooth-form}, rather than the coordinate $s$, defines the
metric across $w=0$.

\emph{Two functions with bounded gradient.}
For a real smooth function $f$, the convention
$g=2\operatorname{Re}\widehat g$ gives
\begin{equation}\label{metric:gradient-convention}
 |\nabla f|_g^2=2|\partial f|_{\widehat g^{-1}}^2.
\end{equation}
Indeed, for a real vector $X$ with $(1,0)$ part $u$,
$g(X,X)=2\widehat g(u,u)$ and
$df(X)=2\operatorname{Re}\partial f(u)$, and the identity follows
by taking dual norms.

Globally on $M_\phi$,
\[
 \partial[-\log(1-t)]
 =\frac{e^\phi\bar w}{1-t}\,\vartheta.
\]
In the coframe of \eqref{metric:smooth-form}, the Hermitian metric is
block diagonal.  Therefore
\[
 |\partial[-\log(1-t)]|_{\widehat g^{-1}}^2
 =\frac{|e^\phi\bar w/(1-t)|^2}{e^\phi/(1-t)^2}=t.
\]
This calculation includes $w=0$, where both sides are zero.  For the
base potential, direct inversion of its complex Hessian gives
\[
 |\partial\psi|_{I^{-1}}^2=|z|^2.
\]
The covector $\partial\psi$ has only base components in the same
coframe, and $G\ge\lambda I$.  Combining these facts with
\eqref{metric:gradient-convention} yields
\begin{equation}\label{metric:gradient-bounds}
 |\nabla[-\log(1-t)]|_g^2=2t<2,
 \qquad
 |\nabla\psi|_g^2\le\frac{2|z|^2}{\lambda}\le\frac2\lambda.
\end{equation}

\emph{Compact containment of metric balls.}
Fix $p\in M_\phi$ and a finite radius $R\ge0$.  Integrating
\eqref{metric:gradient-bounds} along arbitrary piecewise smooth paths
and taking the infimum of their lengths shows that
$d_g(p,q)\le R$ implies
\[
 -\log(1-t(q))\le a_R:=-\log(1-t(p))+\sqrt2R,
 \qquad
 \psi(q)\le b_R:=\psi(p)+\sqrt{2/\lambda}\,R.
\]
This argument does not require a minimizing geodesic.  Set
\[
 \tau_R=1-e^{-a_R}<1,\qquad
 \rho_R=\sqrt{1-e^{-b_R}}<1.
\]
Every such $q=(z,w)$ satisfies $t(q)\le\tau_R$ and
$|z|\le\rho_R$.  Let $m_R$ be the minimum of $\phi$ on the compact
base ball $\{|z|\le\rho_R\}$.  Then
\[
 |w|^2=e^{-\phi(z)}t(q)\le e^{-m_R}\tau_R.
\]
Consequently the metric ball is contained in
\begin{equation}\label{metric:compact-containment}
 K_R=\{(z,w):|z|\le\rho_R,
                  \ e^{\phi(z)}|w|^2\le\tau_R\}.
\end{equation}
This set is closed and bounded in $\C^3$, hence compact, and the
strict bounds $\rho_R<1$ and $\tau_R<1$ place it entirely inside
$M_\phi$.  In particular, a bounded-length curve cannot escape through
the base boundary, through a fiber boundary, or through unbounded fiber
coordinates.  The zero section is an interior locus with the smooth
positive metric \eqref{metric:axis}.

The same compact-containment argument applies to sublevel sets of
$\mathcal K_\lambda$.  Both $\psi$ and $-\log(1-t)$ are nonnegative,
so $\mathcal K_\lambda\le R$ implies
$\psi\le R/\lambda$ and $-\log(1-t)\le R$.  These bounds place the
sublevel set in a compact set of the form
\eqref{metric:compact-containment}; the sublevel set is closed there,
hence compact.  Positivity of its complex Hessian was proved above,
so $\mathcal K_\lambda$ is a strictly plurisubharmonic exhaustion.

\emph{Geodesic completeness.}
Suppose a unit-speed geodesic has a finite maximal forward time $T$.
Its image lies in the compact set $K_T$.  On this set the smooth
positive metric has a uniform positive Euclidean lower bound, so the
geodesic's Euclidean velocity is bounded.  The Christoffel symbols are
also bounded on $K_T$, and the geodesic equation then bounds its
Euclidean acceleration.  Both position and velocity have limits as
time tends to $T$, with limiting position in $K_T\subset M_\phi$.
Local existence for the smooth geodesic equation extends the geodesic
beyond $T$, a contradiction.  The same argument in reverse time proves
geodesic completeness.
\end{proof}

The hypotheses of Proposition~\ref{metric:transfer} imply $h\ge cI>0$,
so Lemma~\ref{metric:complete} applies for every $\lambda>0$ under
consideration.  It remains to choose $\lambda$ large enough for the two
uniform real-sectional bounds.  We do this in the next two sections,
keeping the convention $g=2\operatorname{Re}\widehat g$ throughout.
